\section{Scenario Geometry and Mobility Results}
\label{sec:results_geometry}

The first result block should orient the reader in the simulated environment before any performance metric is discussed. The goal is to show where candidate mounting locations exist, how the selected constellation relates to the urban geometry, and how the user trajectories populate the scene over time.

\begin{figure*}[t]
    \centering
    \begin{subfigure}[t]{0.47\textwidth}
        \centering
        \PaperImage[\linewidth]{\resultroot/scene_layout.png}
        \caption{Top-down Rabot scene layout with candidate sites, selected sites, and simulated mobility overlays.}
        \label{fig:scene_layout}
    \end{subfigure}\hfill
    \begin{subfigure}[t]{0.47\textwidth}
        \centering
        \PaperImage[\linewidth]{\resultroot/trajectory_colormap.png}
        \caption{Trajectory colormap used to visualize the temporal spread of user motion through the walkable space.}
        \label{fig:trajectory_colormap}
    \end{subfigure}
    \caption{Geometry-facing result hooks for the Rabot scenario. These figures should remain the first visual anchor for the simulation setup before performance interpretation begins.}
    \label{fig:geometry_pair}
\end{figure*}

The text accompanying \cref{fig:geometry_pair} should stay factual in this draft. It can describe the modeled environment, the spread of the trajectories, and the low-height infrastructure assumption without yet claiming why one region or path is more favorable than another.

\gilles{Later pass: add 2--3 sentences on how the candidate-wall generation and user-motion spread affect the interpretation of later SINR and schedule figures.}
\FloatBarrier
